\documentclass[10pt]{article}

% ---- NeurIPS style ----
\usepackage{neurips_2026}

% ---- Packages ----
\usepackage[utf8]{inputenc}
\usepackage[T1]{fontenc}
\usepackage{amsmath,amssymb,amsfonts}
\usepackage{graphicx}
\usepackage{booktabs}
\usepackage{hyperref}
\usepackage{natbib}
\usepackage{xcolor}
\usepackage{multirow}
\usepackage{subcaption}
\usepackage{microtype}
\usepackage{enumitem}
\usepackage{wrapfig}

% Allow more floats per page
\renewcommand{\topfraction}{0.85}
\renewcommand{\textfraction}{0.1}
\setcounter{topnumber}{3}

\hypersetup{
  colorlinks=true,
  linkcolor=blue!70!black,
  citecolor=green!50!black,
  urlcolor=blue!70!black
}

% Compact lists
\setlist{nosep,leftmargin=*}

% Warning/collapse marker
\newcommand{\collapse}{\textcolor{red!70!black}{\ding{53}}}
\newcommand{\warn}{{\color{red!70!black}\textbf{!}}}

\title{Where Teachers Teach: The Positional Structure of Distillation Signal}

\author{
  Anonymous Author(s)
}

\begin{document}
\maketitle

% ===========================================================================
% ABSTRACT
% ===========================================================================
\begin{abstract}
On-policy knowledge distillation (KD) trains a student model on reverse KL loss over its own generated sequences.
The standard approach computes loss uniformly across all token positions, implicitly assuming each position contributes equally.
We show this assumption is incorrect and that uniformly weighting positions actively degrades training: teacher signal quality \emph{decays with position}, as later tokens increasingly reflect the teacher confirming the student's prefix rather than providing independent judgment.
We identify an \emph{information-quality paradox}: tokens with the highest KL divergence---seemingly the most informative---produce the worst distillation outcomes (53.4\% on MATH-500), while a contiguous prefix capturing only 46\% of total KL signal achieves 65.85\%.
High-KL tokens correspond to format and style disagreements (LaTeX notation, code blocks), not reasoning quality.
Based on this analysis, we propose \emph{positional distillation}: computing loss only on the first $N$ response tokens---a one-line code change (\texttt{loss\_mask[:, N:] = 0}).
Positional distillation achieves 12--24$\times$ wall-clock speedup and up to 4$\times$ memory reduction, and matches or exceeds full-sequence performance in 4 of 6 math settings (3 families $\times$ \{LoRA, FullFT\}), in all 3 function-calling settings, and across 3 teacher scales (1.7B/4B/8B); the two exceptions are FullFT settings where extended training partially compensates for signal decay (Section~\ref{sec:math_results}).
Cross-model KL profiles motivate a simple descriptive heuristic---a \emph{45\% rule}: in our three families the optimal $N$ falls at the prefix capturing $\sim$40--45\% of cumulative KL signal; we present this as a starting point for tuning $N$, not a validated law. A \emph{cascade effect} additionally propagates early-token improvements to untrained positions, reducing KL by 35--38\% beyond the trained range.
On function calling, distilled students exceed teacher performance by up to 7.3 percentage points, suggesting the method extracts distributional reasoning signal while filtering surface-level noise.
\end{abstract}

% ===========================================================================
% 1. INTRODUCTION
% ===========================================================================
\section{Introduction}
\label{sec:intro}

On-policy knowledge distillation trains a student on reverse KL loss over its own generated sequences, avoiding the train-test mismatch of off-policy methods \citep{agarwal2024onpolicy, gu2024minillm}.
This pipeline is increasingly adopted for distilling reasoning \citep{deepseekr1, chen2025distillingessence} and structured generation capabilities \citep{xu2024survey}.

Standard practice computes loss uniformly over all generated tokens, treating each position as equally informative.
This default has a clean information-theoretic justification---using all tokens maximizes the supervision signal per trajectory---and is therefore the obvious starting point.
Yet in practice, full-sequence on-policy distillation is fragile: across 3 Qwen seeds the mean is 5.2\,pp lower with 1.8$\times$ higher variance than positional, and 1 of 3 seeds \emph{degrades below baseline} after 200 steps of training (seed 123: 50.45\% vs.\ 50.95\% baseline); Gemma drops \emph{below baseline} on every seed, and Llama function calling degrades from 55.3\% to 32.0\%.
These failures are not implementation accidents but symptoms of a structural problem with the assumption that every position carries equal supervision value.

\textbf{Core insight: teacher signal quality decays with position.}
In on-policy distillation, the teacher scores token $t$ conditioned on the student's prefix $y_{<t}$.
At early positions, the teacher provides independent judgment---its distribution reflects genuine knowledge about how to approach a problem.
At later positions, the teacher is increasingly conditioned on the student's prior choices, reducing its role to confirming or rubber-stamping what the student already decided.
We quantify this decay across thousands of on-policy trajectories: mean KL divergence falls from 2.09 at positions 0--50 to 0.43 at positions 150--200 (4.9$\times$ reduction), teacher entropy decreases from 0.26 to 0.17, and teacher-student agreement rises from 78.5\% to 89.5\%.

\textbf{The information-quality paradox.}
A natural response to signal decay would be to train on tokens with the highest KL divergence---targeting where teacher and student disagree most.
We show this intuition is precisely wrong.
Training on the top-100 highest-KL tokens captures 93\% of total KL signal but achieves only 53.4\% on MATH-500.
Training on the first 100 positional tokens captures merely 46\% of KL signal yet achieves 65.85\%.
The resolution: high-KL tokens are dominated by LaTeX formatting disagreements (mean KL = 3.99) and stylistic preferences, not reasoning quality.
The teacher's confusion about format is not useful supervision.

\textbf{Solution: positional distillation.}
We propose \emph{positional distillation}: computing loss only on the first $N$ response tokens and clamping generation accordingly.
This yields a 12$\times$ wall-clock speedup (by generating $N$ tokens instead of $T \gg N$) and 1.9$\times$ memory reduction, while matching or exceeding full-sequence performance across 3 model families, 3 tasks, and 2 training methods.
A \emph{cascade effect} propagates early-token improvements to untrained positions, reducing KL by 35--38\% beyond the trained range.

\textbf{Why not simpler alternatives?} Early stopping leaves the per-token weighting untouched and so still propagates format-disagreement gradients (\S\ref{sec:multiseed}). Magnitude-based token weighting \citep{wang2025beyond8020,huang2025selectkd} is the obvious candidate, but the information-quality paradox (\S\ref{sec:paradox}) shows it would re-weight \emph{toward} the harmful late-position high-KL tokens; we benchmark seven such strategies and all underperform the prefix.

Our contributions are:
\begin{enumerate}
  \item \textbf{Signal-quality analysis} of on-policy KD: teacher supervision quality decays with position across three model families, driven by increasing conditioning on the student's prefix. We propose a \emph{45\% descriptive heuristic} for choosing $N$ (the optimal $N$ in our experiments falls at the prefix capturing 44/40/45\% of cumulative KL for Qwen/Gemma/Llama)---offered as a starting point for tuning, not a validated law.
  \item \textbf{Information-quality paradox.} We demonstrate that high-KL tokens correspond to format/style noise, not reasoning quality. Training on them produces the worst outcome among seven token-selection strategies---contradicting the standard intuition that the largest teacher--student disagreements should be the most informative.
  \item \textbf{Cross-family/task/scale validation.} Positional distillation wins in 4/6 math settings, all 3 function-calling settings, and across 3 teacher scales, with 12--24$\times$ speedup. We also report the two FullFT settings where it underperforms, and analyse why. On function calling, students exceed teacher greedy-decoding accuracy by up to 7.3pp---a known property of soft-label KD that positional distillation amplifies.
  \item \textbf{Cascade effect.} We show that early-token distillation changes entire generation trajectories, with KL reduction extending 35--38\% to untrained positions---the method's benefits are not confined to the trained window.
\end{enumerate}

% Figure 1: Teaser
\begin{figure}[t]
  \centering
  \includegraphics[width=\textwidth]{figures/fig1_teaser.pdf}
  \caption{\textbf{The positional structure of distillation signal.}
  (a) \emph{Information-quality paradox}: top-100 highest-KL tokens (93\% of KL) reach 53.4\%; first-100 positional (46\% of KL) reach 65.85\%. Middle and last-100 fall below baseline.
  (b) \emph{KL decay} across families: Qwen (front-loaded, 2.81$\times$), Gemma, Llama (flat, 1.15$\times$). Dashed lines mark optimal $N$.
  (c) \emph{Cross-family math}: positional matches or exceeds full-seq for all three families; full-seq collapses below baseline on Gemma.
  (d) \emph{Cross-scale teacher}: positional improves the 1.5B student by $+$14.9 to $+$18.0 pp as the Qwen3 teacher scales 1.7B$\to$8B.}
  \label{fig:teaser}
\end{figure}

We take the prefix \emph{as a known empirical recipe} (concurrent work, Section~\ref{sec:concurrent}) and contribute (i)~empirical evidence for \emph{why} it is the right choice in on-policy KD, (ii)~the demonstration that it is \emph{better}, not merely faster, than the full-sequence loss, and (iii)~validation across 3 families, 3 tasks, and 3 teacher scales, including where it underperforms (Section~\ref{sec:limitations}).


% ===========================================================================
% 2. BACKGROUND
% ===========================================================================
\section{Background}
\label{sec:background}

\subsection{On-Policy Knowledge Distillation}

Knowledge distillation transfers knowledge from a teacher $q_\phi$ to a student $p_\theta$ \citep{hinton2015distilling}.
On-policy distillation \citep{agarwal2024onpolicy, gu2024minillm} generates sequences from the student's distribution, with the standard reverse KL objective:
\begin{equation}
\mathcal{L}_\text{full} = \mathbb{E}_{y \sim p_\theta(\cdot|x)} \left[ \sum_{t=1}^{T} D_\mathrm{KL}\!\left( q_\phi(\cdot | x, y_{<t}) \;\|\; p_\theta(\cdot | x, y_{<t}) \right) \right],
\label{eq:fullseq}
\end{equation}
where $x$ is the prompt, $y \sim p_\theta$ is the student-generated response, and $T$ is the response length.
Crucially, the teacher distribution $q_\phi(\cdot|x, y_{<t})$ is conditioned on the student's prefix $y_{<t}$, which becomes increasingly constraining as $t$ grows.

\subsection{Token Selection in Training}

\citet{wang2025beyond8020} show that high-entropy tokens drive RL, finding the top 20\% contribute 80\% of learning signal.
\citet{huang2025selectkd} propose entropy-weighted KD losses; \citet{xie2026adakd} and \citet{kim2026tsdkd} introduce adaptive token-level distillation.
These methods select tokens by signal magnitude, assuming higher-signal tokens are more informative.
We show this assumption fails for on-policy KD: highest-KL tokens are format noise, and \emph{position} matters more than magnitude.

\subsection{Concurrent Work and Our Delta}
\label{sec:concurrent}

\begin{table}[h]
\centering
\caption{\textbf{Concurrent work delta.} Our contribution is the analytical framework, the descriptive 45\% heuristic, and validation breadth---not the prefix heuristic itself.}
\label{tab:concurrent}
\small
\begin{tabular}{@{}lccc@{}}
\toprule
& \textbf{Chen et al.\ 2025} & \textbf{Li et al.\ 2025} & \textbf{This work} \\
\midrule
Setting                          & Off-policy KD       & RL (PPO/GRPO)        & On-policy KD \\
\emph{Why} it works              & Empirical only      & Reward decay (RL)    & Signal decay + paradox \\
Token-selection baselines        & --                  & --                   & 7+ strategies \\
Info.-quality paradox            & --                  & --                   & \checkmark \\
Heuristic for $N$                & --                  & --                   & 45\% rule (descriptive) \\
Family / task / scale validation & 1 / 1 / --          & 1 / 1 / --           & 3 / 3 / 3 \\
\bottomrule
\end{tabular}
\end{table}

\citet{chen2025distillingessence} report prefix truncation as an efficient heuristic for off-policy reasoning distillation, without analyzing why or comparing against token-selection baselines.
\citet{li2025preplan} show reward quality degrades with depth in RL, in a setting different from on-policy KD and without a position-restricted loss.


% ===========================================================================
% 3. TEACHER SIGNAL QUALITY DECAYS WITH POSITION
% ===========================================================================
\section{Teacher Signal Quality Decays with Position}
\label{sec:analysis}

We now ask why the standard loss fails and which tokens actually carry useful supervision. We identify a conditional-scoring problem (\S\ref{sec:conditional}), quantify position-dependent decay (\S\ref{sec:evidence}), uncover the information-quality paradox (\S\ref{sec:paradox}), and show that the resulting KL profile is consistent with the optimal $N$ observed across families (\S\ref{sec:cross_kl}). All analysis uses Qwen2.5-Math-1.5B (student) and Qwen3-1.7B (teacher) with 59,936 on-policy trajectories (mean length 386 tokens) unless noted otherwise.

\subsection{The Conditional Scoring Problem}
\label{sec:conditional}

In on-policy distillation, the teacher scores token $t$ via $q_\phi(\cdot | x, y_{<t}^{\text{student}})$---conditioned on the student's prefix.
Autoregressive conditioning structurally narrows the teacher's high-likelihood continuations as the prefix grows; what is empirically open is \emph{how fast} this narrowing translates into low-information supervision in practice. We frame this as a tractable hypothesis---signal quality decays with position---and the next sections quantify the decay across families and connect it to downstream outcomes.

\subsection{Three Lines of Evidence}
\label{sec:evidence}

\begin{table}[t]
\centering
\caption{\textbf{Teacher signal quality decays with position.} Three complementary metrics on on-policy trajectories show consistent decay. KL divergence drops 4.9$\times$; teacher entropy drops 1.5$\times$; agreement rate (teacher confirms student's top token) rises from 79\% to 90\%.}
\label{tab:decay}
\small
\begin{tabular}{@{}lccccc@{}}
\toprule
\textbf{Metric} & \textbf{Pos 0--50} & \textbf{50--100} & \textbf{100--150} & \textbf{150--200} & \textbf{Ratio (early/late)} \\
\midrule
Mean KL divergence      & 2.09 & 0.94 & 0.53 & 0.43 & 4.9$\times$ \\
Teacher entropy          & 0.26 & 0.24 & 0.18 & 0.17 & 1.5$\times$ \\
Agreement rate (\%)     & 78.5 & 83.7 & 87.6 & 89.5 & --- \\
\bottomrule
\end{tabular}
\end{table}

Three complementary metrics in Table~\ref{tab:decay} converge on the same picture.
\textbf{KL divergence} drops 4.9$\times$ (2.09 $\to$ 0.43); the sharpest decline is in the first 50 tokens, with position 0 (strategy selection) alone at KL = 3.06.
\textbf{Teacher entropy} falls from 0.26 to 0.17 nats: late-position teacher distributions become more peaked, providing less diverse supervision.
\textbf{Agreement rate} (teacher and student top-1 match) rises from 78.5\% to 89.5\%, confirming late-position supervision is largely confirmatory.

\subsection{What Early vs.\ Late Tokens Encode}
\label{sec:token_class}

To understand \emph{what} the teacher disagrees about, we classify 2 million tokens into six categories and compute mean KL by category (Table~\ref{tab:token_class}).

\begin{table}[t]
\centering
\caption{\textbf{Token classification by KL divergence.} The highest-KL category (math\_latex) corresponds to formatting disagreements (\texttt{\textbackslash(}, \texttt{\textbackslash[}), not reasoning quality. Planning tokens (strategy selection) have moderate KL. Math computation tokens (numbers, operators) have the lowest KL.}
\label{tab:token_class}
\small
\begin{tabular}{@{}llrl@{}}
\toprule
\textbf{Category} & \textbf{Example Tokens} & \textbf{Mean KL} & \textbf{Interpretation} \\
\midrule
math\_latex    & \texttt{\textbackslash(}, \texttt{\textbackslash[}, \texttt{\textbackslash\textbackslash} & 3.99 & Format disagreement (noise) \\
planning       & ``To'', ``Therefore'', ``First'' & 1.75 & Strategy selection (signal) \\
structural     & \texttt{**}, \texttt{:}, newlines & 0.89 & Style preference \\
math\_operator & $=$, $+$, $-$                    & 0.38 & Computation (agreement) \\
math\_number   & 0--9                              & 0.28 & Computation (agreement) \\
\bottomrule
\end{tabular}
\end{table}

The highest-KL tokens are LaTeX delimiters (\texttt{\textbackslash(}: KL=10.3; \texttt{\textbackslash[}: KL=13.2 across 28K occurrences)---inline vs.\ display-math conventions that differ between Qwen2.5-Math and Qwen3, a training-data artifact rather than reasoning quality.
Planning tokens (``To'', ``First'', ``Therefore'') have moderate KL and reflect genuine strategic disagreements.
Math computation tokens (numbers, operators) have near-zero KL.

Composition also shifts with position: positions 0--4 are 33\% planning / 2\% math-number, while positions 200--500 are 8\% planning / 18\% math-number.
Early positions encode \emph{strategic} decisions where the teacher has independent judgment; later positions encode \emph{mechanical} execution where agreement is high---explaining the KL decay.

\subsection{The Information-Quality Paradox}
\label{sec:paradox}

Given the token classification, a natural question is: what happens if we train on the highest-KL tokens directly, bypassing position entirely?
We compare seven token selection strategies, each selecting $K=100$ tokens per sequence from full-length generations under a matched training budget (200 steps, 3200 problems, identical optimizer, single seed; only the per-token loss mask differs). Results in Table~\ref{tab:paradox}.

\begin{table}[t]
\centering
\caption{\textbf{The information-quality paradox.} Seven token selection strategies (K=100) on MATH-500. Prefix-100 captures only 46\% of total KL signal but achieves the best performance. Top-KL-100 captures 93\% but performs worst among all methods that train on distillation signal. Middle-100 and Last-100 perform \emph{below the undistilled baseline}, showing that late-position supervision is actively harmful.}
\label{tab:paradox}
\small
\begin{tabular}{@{}lccl@{}}
\toprule
\textbf{Selection Method} & \textbf{KL Coverage (\%)} & \textbf{Best avg@4 (\%)} & \textbf{Trend} \\
\midrule
\textbf{Prefix-100 (ours)}  & \textbf{45.6}  & \textbf{65.85} & Stable \\
Random-100                   & 21.1           & 63.05          & Stable \\
Ent-Teacher-100              & 50.5           & 62.20          & Degrades \\
Full-sequence                & 100.0          & 62.35          & Degrades \\
Ent-Student-100              & 67.0           & 62.70          & Partial \\
Top-KL-100                   & \textbf{93.2}  & \textbf{53.35} & Collapses \\
\midrule
Last-100                     & $\sim$15       & 50.35          & Below baseline \\
Baseline (no distill)        & ---            & 50.95          & --- \\
Middle-100                   & $\sim$30       & 47.80          & Below baseline \\
\bottomrule
\end{tabular}
\end{table}

Top-KL-100 captures 93\% of total KL signal but is 12.5pp worse than Prefix-100 (46\% coverage); Middle-100 and Last-100 fall \emph{below} the undistilled baseline. More signal concentration yields worse outcomes.

\textbf{Resolution.} High-KL tokens are dominated by LaTeX formatting disagreements (\S\ref{sec:token_class})---the teacher is confused by the student's notation choice, not superior to its reasoning, so training on these positions teaches format preferences, not mathematics. The prefix succeeds for the opposite reason: in autoregressive generation, early decisions condition all subsequent tokens, so improving them rewrites the whole trajectory; scattered high-KL corrections do not propagate.

\subsection{Cross-Model KL Profiles}
\label{sec:cross_kl}

The positional structure of KL divergence varies across model families (Table~\ref{tab:kl_profiles}).

\begin{table}[t]
\centering
\caption{\textbf{Cross-model KL profiles---the 45\% descriptive heuristic.} Front-loaded KL profiles (Qwen) need fewer tokens; flat profiles (Llama) need more. In all three families the optimal $N$ falls at the prefix capturing $\sim$40--45\% cumulative KL signal. We propose this as a starting heuristic for tuning $N$ on a new model family, not a validated law---validation on additional families is left to future work.}
\label{tab:kl_profiles}
\small
\begin{tabular}{@{}lcccc@{}}
\toprule
\textbf{Student $\to$ Teacher} & \textbf{Ratio} & \textbf{Best $N$} & \textbf{Cumul.\ KL} & \textbf{Shape} \\
\midrule
Qwen M-1.5B $\to$ Q3-1.7B & 2.81$\times$ & 100 & $\sim$44\% & Front-loaded \\
Gemma 2B $\to$ 3-4B        & $\sim$2.5$\times$ & 50--100 & $\sim$40\% & Front-loaded \\
Llama 1B $\to$ 3.1-8B      & 1.15$\times$ & 200 & $\sim$45\% & Flat \\
\bottomrule
\end{tabular}
\end{table}

Qwen is strongly front-loaded (ratio 2.81$\times$, optimum $N=100$); Llama is nearly flat (1.15$\times$, KL $\sim$0.18 everywhere, optimum $N=200$).
Despite these shape differences, the three optimal-$N$ values fall in a narrow $\sim$40--45\% cumulative-KL band (44/40/45\% for Qwen/Gemma/Llama).
We treat this as a \emph{descriptive heuristic} for choosing an initial $N$, not a predictive law: at $<$40\% coverage the prefix is too short (Llama $N{=}150$, 34\%: 18.75 vs.\ 22.45 at $N{=}200$); at $>$50\% the loss starts incorporating the late-position high-KL tokens shown to be harmful in \S\ref{sec:paradox}.
Whether the band reflects a deep property of teacher signal-to-noise structure is an open question (\S\ref{sec:limitations}).


% ===========================================================================
% 4. METHOD
% ===========================================================================
\section{Method: Positional Distillation}
\label{sec:method}

The analysis pinpoints \emph{where} useful supervision lives---the early-position prefix---so the method follows directly: restrict the standard loss to that prefix and clamp generation accordingly.

\subsection{Formulation}

Positional distillation restricts the loss in Eq.~\ref{eq:fullseq} to the first $N$ response tokens:
\begin{equation}
\mathcal{L}_{\text{pos-}N} = \mathbb{E}_{y \sim p_\theta(\cdot|x)} \left[ \sum_{t=1}^{N} D_\mathrm{KL}\!\left( q_\phi(\cdot | x, y_{<t}) \;\|\; p_\theta(\cdot | x, y_{<t}) \right) \right].
\label{eq:pos}
\end{equation}

The implementation requires two changes: a one-line loss mask (\texttt{loss\_mask[:, N:] = 0}) and clamping generation length (\texttt{max\_new\_tokens = N}) so the student generates exactly $N$ tokens per problem.
The teacher scores these $N$ tokens with a single forward pass.
No changes to the model architecture, optimizer, or training loop are required.

\subsection{Efficiency}

\begin{table}[t]
\centering
\caption{\textbf{Efficiency comparison.} Per-step timing and memory on a single A100 40GB GPU with Qwen2.5-Math-1.5B student and Qwen3-1.7B teacher, batch size 16. Positional distillation achieves 12--24$\times$ speedup and up to 4$\times$ memory reduction.}
\label{tab:efficiency}
\small
\begin{tabular}{@{}lccccc@{}}
\toprule
\textbf{Component} & \textbf{Full-seq} & \textbf{Pos-100} & \textbf{Speedup} \\
\midrule
Generation               & 122--180s & 6s  & 20--36$\times$ \\
Teacher scoring           & 5--7s     & 1s  & 5--7$\times$ \\
Training (fwd+bwd)       & 6--7s     & 3.7s & 1.7--1.9$\times$ \\
\midrule
\textbf{Total per step}  & \textbf{133--194s} & \textbf{8--11s} & \textbf{12--24$\times$} \\
\textbf{200 steps}       & \textbf{7.4--10.8h} & \textbf{27--37min} & \\
\textbf{Peak GPU memory} & \textbf{17--39.5 GB} & \textbf{9--9.6 GB} & \textbf{1.9--4$\times$} \\
\bottomrule
\end{tabular}
\end{table}

The speedup (Table~\ref{tab:efficiency}) comes overwhelmingly from generation: producing 100 tokens instead of 2048--3584 reduces it from 122--180s to 6s ($>$95\% of total savings). Memory savings come from shorter training-time activations (100 vs.\ 1300+ tokens). At 9.6\,GB peak, positional distillation enables single-GPU training where full-sequence requires multi-GPU or offloading.

\subsection{Selecting $N$}
\label{sec:selecting_n}

We recommend the following heuristic: generate $\sim$1000 on-policy trajectories, compute per-position KL, and select $N$ at approximately 40--50\% of cumulative KL coverage.
For front-loaded profiles (Qwen-like), this yields $N = 50$--100.
For flat profiles (Llama-like), $N = 150$--200.
In practice, performance is robust to the exact choice: for Qwen, $N \in [50, 200]$ all achieve within 1pp of the optimum (Table~\ref{tab:position_sweep}).


% ===========================================================================
% 5. EXPERIMENTS
% ===========================================================================
\section{Experiments}
\label{sec:experiments}

We validate three claims: positional distillation matches or beats full-sequence in realistic settings (\S\ref{sec:math_results}, \S\ref{sec:funcall_results}); the 45\% heuristic generalizes beyond Qwen math (\S\ref{sec:cross_kl}, \S\ref{sec:math_results}); and the structural benefits (multi-seed stability, cascade) are not checkpoint-selection artifacts (\S\ref{sec:multiseed}, \S\ref{sec:ablations}).

\subsection{Setup}
\label{sec:setup}

\textbf{Model families.}
We evaluate three student-teacher pairs spanning different architectures and capability gaps:
\begin{itemize}
  \item \textbf{Qwen}: Qwen2.5-Math-1.5B $\to$ Qwen3-1.7B (cross-architecture, math-specialized student)
  \item \textbf{Gemma}: Gemma-2-2B-IT $\to$ Gemma-3-4B-IT (cross-generation, instruction-tuned)
  \item \textbf{Llama}: Llama-3.2-1B-Instruct $\to$ Llama-3.1-8B-Instruct (cross-generation, 8$\times$ size gap)
\end{itemize}
Additionally, we test Qwen2.5-Math-1.5B $\to$ Qwen3-4B (larger teacher) and Qwen3-4B $\to$ Qwen3-8B (same-family negative result).

\textbf{Tasks.}
(i)~\textbf{MATH-500} \citep{hendrycks2021measuring}: mathematical reasoning, evaluated with avg@4 (4 samples, temperature 0.7).
(ii)~\textbf{BFCL}: function calling (600 problems), evaluated with full accuracy (function name + arguments).
(iii)~\textbf{HumanEval/MBPP} \citep{chen2021evaluating, austin2021program}: code generation, evaluated with pass@1 (temperature 0.0).

\textbf{Training.}
LoRA \citep{hu2022lora}: $r=32$, $\alpha=64$, all linear layers, $\text{lr}=5\times10^{-5}$.
Full fine-tuning: $\text{lr}=5\times10^{-6}$.
All: 200 steps, 3200 training problems, $n_{\text{samples}}=1$, batch size 16, cosine LR schedule, checkpoints every 50 steps.
We report best checkpoint across steps 50--200.

\subsection{Main Results: Mathematical Reasoning}
\label{sec:math_results}

\begin{table}[t]
\centering
\caption{\textbf{Math results (MATH-500 avg@4).} Positional distillation wins in 4 of 6 settings. Gemma full-sequence drops \emph{below baseline} in both LoRA and FullFT. LoRA outperforms FullFT by 8--13pp across all families. The exceptions are \emph{Qwen FullFT} (extended 400-step training lets full-sequence catch up) and \emph{Llama FullFT} (flat KL profile makes the prefix less informative)---we discuss when positional underperforms in Section~\ref{sec:discussion}.}
\label{tab:math_main}
\small
\begin{tabular}{@{}lcccccc@{}}
\toprule
& \multicolumn{2}{c}{\textbf{Qwen}} & \multicolumn{2}{c}{\textbf{Gemma}} & \multicolumn{2}{c}{\textbf{Llama}} \\
\cmidrule(lr){2-3} \cmidrule(lr){4-5} \cmidrule(lr){6-7}
& LoRA & FullFT & LoRA & FullFT & LoRA & FullFT \\
\midrule
Baseline                & 50.95 & 50.95 & 13.45 & 13.45 & 15.20 & 15.20 \\
\midrule
Best positional         & \textbf{66.65} & 54.55 & \textbf{25.80} & \textbf{26.70} & \textbf{22.45} & 17.80 \\
\quad ($N$)             & \footnotesize{($N$=50)} & \footnotesize{($N$=150)} & \footnotesize{($N$=50)} & \footnotesize{($N$=100)} & \footnotesize{($N$=200)} & \footnotesize{($N$=300)} \\
Full-sequence           & 62.35 & \textbf{57.10} & 11.70{\color{red}\textbf{*}} & 13.90{\color{red}\textbf{*}} & 20.65 & \textbf{20.10} \\
\midrule
$\Delta$ (pos $-$ full) & +4.30 & $-$2.55 & +14.10 & +12.80 & +1.80 & $-$2.30 \\
\bottomrule
\end{tabular}
\\[2pt]
{\footnotesize {\color{red}\textbf{*}} Below baseline. Qwen LoRA reports best checkpoint (pos-50 step 150); FullFT reports extended training to 400 steps.}
\end{table}

\textbf{Positional distillation wins in 4 of 6 settings} (Table~\ref{tab:math_main}). The two losses are both FullFT: Qwen FullFT recovers from signal decay via extended 400-step training, and Llama FullFT has a nearly-flat KL profile (1.15$\times$) that offers no informative prefix to isolate. We conjecture LoRA's faster effective updates compound late-position noise into instability while FullFT's slower updates absorb it (direct gradient-magnitude evidence is left to future work).
\textbf{Gemma full-sequence drops below baseline} (11.70/13.90 vs.\ 13.45), while positional ($N{=}50$) nearly doubles accuracy to 25.80.
\textbf{LoRA outperforms FullFT by 8--13pp} across families.
\textbf{Optimal $N$ tracks KL profile}: front-loaded Qwen peaks at $N{=}50$--100; flat-profile Llama needs $N{=}200$--300, matching the heuristic in \S\ref{sec:cross_kl}.

\subsection{Main Results: Function Calling}
\label{sec:funcall_results}

\begin{table}[t]
\centering
\caption{\textbf{Function calling results (BFCL full accuracy).} Students exceed teacher accuracy on Qwen (+7.3pp) and Gemma (+5.9pp) despite the teacher not generating valid JSON output. Llama full-sequence collapses from 55.3\% baseline to 32.0\%. Positional distillation outperforms full-sequence across all three families.}
\label{tab:funcall}
\small
\begin{tabular}{@{}lccc@{}}
\toprule
& \textbf{Qwen} & \textbf{Gemma} & \textbf{Llama} \\
\midrule
Teacher accuracy & 54.0\% & 25.0\% & 63.7\% \\
Student baseline & 2.7\% & 0\% & 55.3\% \\
\midrule
Best positional  & \textbf{61.3\%} ($N$=100) & \textbf{30.9\%} ($N$=50) & \textbf{59.0\%} ($N$=150) \\
\quad vs.\ teacher & {\color{green!50!black}+7.3pp} & {\color{green!50!black}+5.9pp} & $-$4.7pp \\
Full-sequence    & 58.2\% & 3.9\% & 32.0\% \\
\bottomrule
\end{tabular}
\end{table}

Two patterns stand out (Table~\ref{tab:funcall}).
First, \textbf{students exceed teachers} on Qwen (+7.3pp) and Gemma (+5.9pp) despite neither teacher generating valid JSON (Qwen outputs natural language, Gemma outputs Python-style \texttt{tool\_code}). The teacher's distribution assigns non-trivial mass to correct function names/arguments even when its greedy output does not; positional distillation amplifies this by focusing on the $<$100-token window where function-call structure resides.
Second, \textbf{Llama full-sequence collapses} from 55.3\% to 32.0\% ($-$23.3pp); positional ($N{=}150$) preserves baseline at 59.0\%.

\subsection{Multi-Seed Stability}
\label{sec:multiseed}

\begin{table}[t]
\centering
\caption{\textbf{Multi-seed stability (Qwen 1.5B$\to$1.7B, n1bs16 LoRA, MATH-500 avg@4 at step 200, last-boxed extraction).} Positional distillation is more stable across seeds (std 2.96 vs.\ 5.40, $1.8\times$ tighter) and has a higher mean (+5.2pp). Pos-100 wins on every seed; the worst full-sequence seed (123) \emph{degrades below baseline} (50.45\% vs.\ 50.95\% no-distillation baseline) -- 1 of 3 full-sequence runs actively makes the model worse.}
\label{tab:multiseed}
\small
\begin{tabular}{@{}lccccc@{}}
\toprule
\textbf{Method} & \textbf{Seed 42} & \textbf{Seed 123} & \textbf{Seed 456} & \textbf{Mean $\pm$ Std} & \textbf{$\Delta$ baseline} \\
\midrule
Pos-100 (LoRA) & \textbf{65.85} & 61.00 & 60.40 & $\mathbf{62.42 \pm 2.96}$ & $+11.5 \pm 2.96$ \\
Full-seq (LoRA) & 63.55 & \textbf{50.45} & 57.55 & $57.18 \pm 5.40$ & $+6.2 \pm 5.40$ \\
\bottomrule
\end{tabular}
\\[2pt]
{\footnotesize All values are MATH-500 \texttt{avg@4} at step 200 with $n_{\text{samples}}{=}4$, $T{=}0.7$, last-boxed answer extraction (the standard convention). Baseline (no distillation): 50.95\%.}
\end{table}

On three Qwen seeds (Table~\ref{tab:multiseed}), full-sequence is both lower and noisier: mean 57.18\% (std 5.40) vs.\ pos-100's 62.42\% (std 2.96). The worst full-sequence seed (123) reaches only 50.45\%, \emph{below} the no-distillation baseline of 50.95\% -- 200 steps of full-sequence training made the model worse than no training at all (a $-0.5$pp degradation on this seed). Pos-100 instead clears baseline by at least $+9.5$pp on every seed. Inspecting outputs of the underperforming full-sequence seeds reveals the failure mode: the model has learned to repeat \texttt{\textbackslash boxed\{\}} dozens of times per response, a degenerate reward-hack pattern that the standard last-boxed answer extraction parses as the model's final answer. We report this under the canonical metric without compensating extraction strategies; the extraction failure \emph{is} the failure.

\subsection{Ablations}
\label{sec:ablations}

\begin{table}[t]
\centering
\caption{\textbf{Position sweep (Qwen LoRA, MATH-500 avg@4).} Performance is robust across $N \in [50, 200]$, with all values achieving within 1pp of the optimum. Both very small ($N=5$) and very large ($N=$ full-seq) positions underperform.}
\label{tab:position_sweep}
\small
\begin{tabular}{@{}lcccccccc@{}}
\toprule
$N$ & 5 & 10 & 20 & 50 & 100 & 150 & 200 & Full \\
\midrule
avg@4 & 56.50 & 59.50 & 60.20 & \textbf{66.65} & 65.85 & 66.65 & 66.05 & 62.35 \\
\bottomrule
\end{tabular}
\\[2pt]
{\footnotesize $N=5$--$50$ from small-scale (200 problems, $n$=16); $N=50$--Full from main config (3200 problems, $n$=1).}
\end{table}

\textbf{Position sweep.}
On Qwen, $N \in [50, 200]$ all achieve within 1pp of the optimum (66.65; Table~\ref{tab:position_sweep}); $N{=}5$ underfits (56.50) and full-sequence (62.35) is diluted by late-position noise. The broad optimum makes $N$ robust.

\textbf{Teacher scaling.}
Scaling the teacher from Qwen3-1.7B to 4B improves from 65.85\% to 68.95\% (+3.1pp), but same-family distillation (Qwen3-4B $\to$ 8B) yields $-$0.9pp, confirming diminishing returns near capacity ceilings.

\textbf{Cascade effect.}
We measure mean teacher--student KL per position bin on 5K held-out trajectories sampled from the raw and pos-200-distilled (step 100) students. Pos-200 distillation reduces KL by 35--38\% at \emph{untrained} positions 200--400 (Table~\ref{tab:cascade}), confirming early decisions propagate through autoregressive generation.

\begin{table}[t]
\centering
\caption{\textbf{Cascade effect.} Pos-200 distillation reduces KL even at untrained positions (200--400), by 35--38\%. Full-sequence shows larger tail reductions but also exhibits repetition artifacts.}
\label{tab:cascade}
\small
\begin{tabular}{@{}lcccc@{}}
\toprule
\textbf{Position Range} & \textbf{Raw} & \textbf{Pos-200 s100} & \textbf{$\Delta$} & \textbf{Trained?} \\
\midrule
0--50    & 2.064 & 1.015 & $-$51\% & Yes \\
50--100  & 0.759 & 0.321 & $-$58\% & Yes \\
100--200 & 0.441 & 0.237 & $-$46\% & Yes \\
200--300 & 0.382 & 0.238 & $-$38\% & \textbf{No} \\
300--400 & 0.331 & 0.216 & $-$35\% & \textbf{No} \\
\bottomrule
\end{tabular}
\end{table}

\subsection{Degradation Analysis}
\label{sec:degradation}

Full-sequence distillation exhibits three failure modes:
(1)~\textbf{Repetition reward-hacking (Qwen)}: under last-boxed extraction (the canonical metric), full-sequence scores 57.18\,$\pm$\,5.40 across seeds versus positional 62.42\,$\pm$\,2.96; one of three seeds (seed 123) \emph{degrades to 50.45\%, below the 50.95\% no-distillation baseline}---i.e., 200 steps of full-sequence training made the model worse than no training at all. Inspection shows responses repeating \texttt{\textbackslash boxed\{\}} dozens of times---a degenerate optimum that satisfies the on-policy reverse-KL objective but is penalized by the canonical metric.
(2)~\textbf{Below-baseline collapse (Gemma)}: full-sequence achieves 11.70\% LoRA / 13.90\% FullFT, both below the 13.45\% baseline.
(3)~\textbf{Capability destruction (Llama funcall)}: function calling drops from 55.3\% to 32.0\% despite 99\%+ parse rate---full-sequence training on 2048 tokens overwhelms the 33-token function call structure.


% ===========================================================================
% 6. DISCUSSION
% ===========================================================================
\section{Discussion}
\label{sec:discussion}

\subsection{Why Positional Distillation Works}

Three mechanisms compose: (1)~\textbf{Signal-quality filtering}---early positions have lower agreement (75\%) and higher teacher entropy (0.26) than late positions (89\%, 0.17), so restricting to the prefix avoids diluting high-quality signal with confirmation; (2)~\textbf{Noise avoidance}---late-position KL is dominated by format disagreements whose high-variance gradients cause repetition collapse and below-baseline degradation; (3)~\textbf{Cascade}---improving early strategy/structure decisions rewrites entire trajectories, evidenced by the 35--38\% KL reduction at untrained positions.

\subsection{Why Students Exceed Teachers}

On function calling, students exceed teacher greedy-decoding accuracy by up to 7.3pp because on-policy distillation decouples the teacher's \emph{distributional knowledge} from its \emph{generation ability}: the Qwen3-1.7B teacher's greedy output is natural language, yet its distribution still assigns non-trivial mass to correct function names and arguments, which the on-policy student learns to emit. Positional distillation amplifies this by focusing on the $<$100-token window where structure lives, avoiding positions where the teacher reverts to prose. The implication is general: a model's soft distribution can encode capabilities its argmax never expresses.

\subsection{Limitations}
\label{sec:limitations}

\textbf{Where positional underperforms.} In 2 of 6 main math settings (Qwen FullFT, Llama FullFT) full-sequence matches or slightly beats positional. Both are FullFT (slower updates absorb late-position noise) and involve relatively flat KL profiles; practitioners running long FullFT schedules on flat-profile teachers should not assume positional dominates.
\textbf{Generality of the 45\% heuristic.} Calibrated on three families ($\le$8B); we present it as a starting point, not a law (\S\ref{sec:cross_kl}).
\textbf{Mechanism evidence.} The \emph{why} analysis (signal decay, token classification, paradox) is on Qwen; we infer the same mechanism for Gemma/Llama from consistent KL profiles, but direct token-level analysis on those families is future work.
\textbf{Position vs.\ content confound.} The strongest critical reading of our results is that what the prefix really filters is \emph{format/style noise} (LaTeX, code blocks), not late-position tokens per se. The cleanest disentangling experiment is to run \emph{full-sequence loss with format tokens masked out}: if it matches positional, the operative variable is content, not position. We did not run this ablation (it requires per-token classifiers across all positions on all three families) and we flag it as the highest-value follow-up. The cross-family KL-decay evidence (\S\ref{sec:evidence}) and cascade effect (\S\ref{sec:cascade}) suggest position carries independent signal, but our current evidence does not isolate the two factors.
\textbf{Task and scale scope.} Validation covers structured reasoning (math, function calling, code) on students $\le$4B and teachers $\le$8B; benefits for open-ended generation and at larger scale are open. We evaluate reverse KL only.


% ===========================================================================
% 7. CONCLUSION
% ===========================================================================
\section{Conclusion}

On-policy knowledge distillation has a positional structure: teacher signal quality decays with position as supervision shifts from independent judgment to passive confirmation, and this decay drives an information-quality paradox in which the highest-KL tokens are dominated by format noise rather than reasoning signal. Together these observations motivate positional distillation (\texttt{loss\_mask[:, N:] = 0}) not as an efficiency hack, but as a principled response to where the teacher actually carries information. Across 3 families, 3 tasks, and 2 training methods, it matches or exceeds full-sequence in 4 of 6 math settings and all function-calling settings with 12--24$\times$ speedup; across three Qwen seeds the mean is 5.2\,pp higher and variance is 1.8$\times$ tighter (std 5.40 $\to$ 2.96), and 1 of 3 full-sequence seeds degrades below baseline whereas every positional seed clears it by $\ge$\,9\,pp. On function calling, distilled students exceed teacher greedy-decoding accuracy by up to 7.3pp, extracting distributional knowledge the teacher cannot itself emit. The dependence of optimal $N$ on the teacher's KL profile points to a single broader claim: in on-policy distillation, \emph{where} the teacher supervises matters at least as much as \emph{how much}.


% ===========================================================================
% REFERENCES
% ===========================================================================
\bibliographystyle{plainnat}
\bibliography{references_extended}


% ===========================================================================
% APPENDIX
% ===========================================================================
\newpage
\appendix

\section{Full Per-Step Results}
\label{app:perstep}

\subsection{Qwen LoRA Math (MATH-500, avg@4/maj@4/pass@4)}

Baseline (undistilled Qwen2.5-Math-1.5B): 50.95\% avg@4, 61.2\% maj@4, 72.8\% pass@4.

\begin{table}[h]
\centering
\caption{Qwen LoRA per-step math results (3200 problems, $n$=1, bs=16).}
\footnotesize
\begin{tabular}{@{}lcccccccccccc@{}}
\toprule
& \multicolumn{3}{c}{\textbf{Step 50}} & \multicolumn{3}{c}{\textbf{Step 100}} & \multicolumn{3}{c}{\textbf{Step 150}} & \multicolumn{3}{c}{\textbf{Step 200}} \\
\cmidrule(lr){2-4} \cmidrule(lr){5-7} \cmidrule(lr){8-10} \cmidrule(lr){11-13}
& avg & maj & pass & avg & maj & pass & avg & maj & pass & avg & maj & pass \\
\midrule
Pos-50  & 62.4 & 69.4 & 77.2 & 66.1 & 72.0 & 79.4 & \textbf{66.7} & 71.0 & 81.0 & 64.9 & 71.2 & 79.6 \\
Pos-100 & 63.8 & 70.0 & 79.8 & 64.5 & 68.4 & 78.4 & 65.2 & 69.6 & 80.2 & \textbf{65.9} & 70.8 & 79.8 \\
Pos-150 & 65.4 & 66.8 & 79.0 & \textbf{66.7} & 67.0 & 81.0 & 65.3 & 66.3 & 78.2 & 65.8 & 67.3 & 80.0 \\
Pos-200 & \textbf{66.1} & 71.2 & 81.0 & 64.7 & 68.4 & 79.8 & 65.1 & 70.0 & 80.6 & 65.6 & 71.2 & 80.6 \\
Full    & 61.0 & 66.8 & 74.6 & 62.0 & 68.4 & 75.2 & \textbf{62.4} & 69.4 & 74.6 & 61.2 & 65.2 & 75.0 \\
\bottomrule
\end{tabular}
\end{table}

\subsection{Qwen FullFT Math (MATH-500, avg@4)}

Baseline: 50.95\%.

\begin{table}[h]
\centering
\caption{Qwen FullFT per-step math results (extended to 400 steps, $n$=1, bs=16).}
\small
\begin{tabular}{@{}lccccccccc@{}}
\toprule
\textbf{Step} & 50 & 100 & 150 & 200 & 250 & 300 & 350 & 400 & \textbf{Best} \\
\midrule
Pos-100  & 50.85 & 52.35 & 53.00 & \textbf{53.15} & 49.20 & 49.90 & 50.75 & 49.95 & 53.15 \\
Pos-150  & 51.40 & 53.10 & 53.70 & 53.25 & 53.30 & \textbf{54.55} & 53.85 & 52.85 & 54.55 \\
Pos-200  & 51.30 & \textbf{54.05} & 53.40 & 53.75 & 52.55 & 53.75 & 53.10 & 53.85 & 54.05 \\
Pos-300  & 50.55 & \textbf{53.70} & 53.95 & 53.00 & 54.45 & \textbf{55.10} & 54.95 & 54.25 & 55.10 \\
Full-seq & 51.40 & 53.55 & 54.95 & 55.95 & \textbf{57.05} & \textbf{57.10} & 56.05 & 56.20 & \textbf{57.10} \\
\bottomrule
\end{tabular}
\end{table}

\subsection{Gemma LoRA Math (MATH-500, avg@4)}

Baseline (gemma-2-2b-it): 13.45\%.

\begin{table}[h]
\centering
\caption{Gemma LoRA per-step math results.}
\small
\begin{tabular}{@{}lcccc@{}}
\toprule
\textbf{Step} & 50 & 100 & 150 & 200 \\
\midrule
Pos-100 (LoRA) & \textbf{27.20} & 26.95 & 26.60 & 27.45 \\
Fullseq (LoRA) & --- & --- & --- & 11.70 \\
\midrule
Pos-100 (FullFT) & 25.20 & \textbf{26.70} & 25.45 & 26.65 \\
Fullseq (FullFT) & 11.15 & 9.50 & 13.00 & \textbf{13.90} \\
\bottomrule
\end{tabular}
\end{table}

\subsection{Llama LoRA/FullFT Math (MATH-500, avg@4)}

Baseline (Llama-3.2-1B-Instruct): 15.20\%.

\begin{table}[h]
\centering
\caption{Llama per-step math results.}
\small
\begin{tabular}{@{}lccl@{}}
\toprule
\textbf{Method} & \textbf{Config} & \textbf{Best avg@4} & \textbf{Notes} \\
\midrule
LoRA pos-100  & & 12.35 & Below baseline \\
LoRA pos-150  & & 18.75 & \\
LoRA pos-200  & & \textbf{22.45} & Best LoRA \\
LoRA pos-300  & & 20.75 & \\
LoRA fullseq  & & 20.65 & \\
\midrule
FullFT pos-100 & & 12.80 & Below baseline \\
FullFT pos-200 & & 16.80 & \\
FullFT pos-300 & & 17.80 & \\
FullFT fullseq & & \textbf{20.10} & Best FullFT \\
\bottomrule
\end{tabular}
\end{table}


\section{Token Classification Details}
\label{app:token_class}

Token classification was performed on 2,015,036 tokens from 10,000 trajectories using rule-based categorization over the Qwen2.5-Math-1.5B tokenizer vocabulary.

\textbf{Categories:}
\begin{itemize}
\item \textbf{planning}: Sentence-initial strategy words (``To'', ``Let'', ``First'', ``Therefore'', ``Thus'', ``Since'', ``Given'', etc.)
\item \textbf{structural}: Formatting tokens (\texttt{**}, \texttt{:}, newlines, \texttt{\$}, parentheses, brackets)
\item \textbf{math\_latex}: LaTeX-specific commands (\texttt{\textbackslash(}, \texttt{\textbackslash[}, \texttt{\textbackslash)}, \texttt{\textbackslash]}, \texttt{\textbackslash frac}, etc.)
\item \textbf{math\_number}: Digit tokens (0--9)
\item \textbf{math\_operator}: Arithmetic operators ($=$, $+$, $-$, $\times$, $/$ etc.)
\item \textbf{continuation}: All remaining tokens (common words, word fragments, etc.)
\end{itemize}

\textbf{Position 0 deep dive.}
Position 0 has mean KL = 7.71 (10$\times$ the overall mean of 1.01), driven by strategy selection.
The dominant token is ``To'' (75.1\% of trajectories, KL = 7.32), as in ``To find/determine/solve...''.
The teacher prefers more diverse openings: ``First'' (KL = 14.0), ``Step'' (KL = 14.6), ``Solution'' (KL = 21.9).
This single position accounts for 3.8\% of total KL signal across the entire trajectory.

\textbf{Category composition by position range:}

\begin{table}[h]
\centering
\small
\begin{tabular}{@{}lcccccc@{}}
\toprule
\textbf{Category} & \textbf{0--4} & \textbf{5--19} & \textbf{20--49} & \textbf{50--99} & \textbf{100--199} & \textbf{200--499} \\
\midrule
planning (\%)        & 32.8 & 15.9 & 12.2 & 9.7 & 8.8 & 7.9 \\
structural (\%)      & 7.8  & 19.4 & 25.0 & 30.8 & 34.8 & 39.3 \\
math\_number (\%)    & 1.7  & 9.5  & 11.3 & 15.2 & 17.3 & 18.1 \\
math\_operator (\%)  & 0.3  & 2.5  & 3.3  & 5.3  & 5.8  & 4.5 \\
math\_latex (\%)     & 1.4  & 3.1  & 2.7  & 3.8  & 4.8  & 5.9 \\
continuation (\%)    & 56.1 & 49.6 & 45.5 & 35.1 & 28.5 & 24.4 \\
\bottomrule
\end{tabular}
\end{table}


\section{Code Generation Results}
\label{app:coding}

\begin{table}[h]
\centering
\caption{Coding results (HumanEval/MBPP pass@1). Qwen LoRA only. Baseline: HE 32.9\%, HE+ 29.3\%.}
\small
\begin{tabular}{@{}lccccl@{}}
\toprule
\textbf{Method} & \textbf{Best HE} & \textbf{Best HE+} & \textbf{Best MBPP} & \textbf{Best MBPP+} & \textbf{Step} \\
\midrule
LoRA pos-50   & \textbf{42.1} & \textbf{36.6} & 46.6 & 41.5 & 350 \\
LoRA pos-100  & \textbf{42.1} & 36.0          & 49.2 & 44.4 & 150 \\
LoRA pos-150  & 41.5          & \textbf{37.8} & 49.5 & 44.7 & 250 \\
LoRA fullseq  & 40.2          & 35.4          & \textbf{52.6} & \textbf{46.3} & 50 \\
\midrule
FullFT pos-250 & \textbf{37.8} & \textbf{31.7} & \textbf{54.2} & \textbf{47.4} & 350 \\
FullFT pos-150 & 37.2          & 31.1          & 54.2 & 46.8 & 300/400 \\
FullFT fullseq & 36.6          & 31.7          & 53.2 & 46.0 & 150 \\
\bottomrule
\end{tabular}
\end{table}

\textbf{Key findings:}
\begin{itemize}
\item LoRA positional variants achieve the best HumanEval scores (42.1\% for pos-50 and pos-100), while full-sequence achieves the best MBPP scores (52.6\%).
\item LoRA full-sequence HumanEval degrades from 40.2\% (step 50) to 26.8\% (step 400), a 13.4pp drop. Positional variants show no such degradation.
\item FullFT results are more stable but uniformly lower than LoRA, consistent with math findings.
\item The HumanEval/MBPP trade-off may reflect that shorter prefixes encourage concise solutions (beneficial for HumanEval's simple functions) while longer training encourages more complete implementations (beneficial for MBPP's broader scope).
\end{itemize}


\section{Additional Student-Teacher Pairs}
\label{app:additional_pairs}

\subsection{Teacher Scaling: Qwen2.5-Math-1.5B with Larger Teachers}

\begin{table}[h]
\centering
\caption{Teacher scaling results (LoRA pos-100, MATH-500 avg@4).}
\small
\begin{tabular}{@{}llccc@{}}
\toprule
\textbf{Student} & \textbf{Teacher} & \textbf{Baseline} & \textbf{Best avg@4} & \textbf{$\Delta$} \\
\midrule
M-1.5B & Qwen3-1.7B & 50.95 & 65.85 & +14.9 \\
M-1.5B & Qwen3-4B   & 50.95 & \textbf{68.95} & \textbf{+18.0} \\
M-1.5B & Qwen3-8B   & 50.95 & 67.85 & +16.9 \\
\bottomrule
\end{tabular}
\end{table}

Scaling the teacher from 1.7B to 4B yields a 3.1pp improvement.
The 8B teacher underperforms the 4B teacher, possibly because the larger capacity gap makes the teacher's distribution harder for the small student to learn from.

\subsection{Same-Family Negative Results}

\begin{table}[h]
\centering
\caption{Same-family distillation (Qwen3 family, LoRA pos-100, MATH-500 avg@4).}
\small
\begin{tabular}{@{}llccl@{}}
\toprule
\textbf{Student} & \textbf{Teacher} & \textbf{Baseline} & \textbf{Best avg@4} & \textbf{$\Delta$} \\
\midrule
Qwen3-1.7B & Qwen3-4B & 69.20 & 69.20 & $\pm$0.0 \\
Qwen3-1.7B & Qwen3-8B & 69.20 & 69.15 & $-$0.05 \\
Qwen3-4B   & Qwen3-8B & 77.95 & 77.05 & $-$0.9 \\
\bottomrule
\end{tabular}
\end{table}

Same-family distillation within the Qwen3 family shows zero or negative gains on math.
These students are already near their capacity ceiling: Qwen3-1.7B achieves 69.2\% (vs.\ teacher's 69.2\% at temperature 0.7), and Qwen3-4B achieves 78.0\% (vs.\ teacher's 77.7\%).
The teacher provides no additional mathematical reasoning signal when the student is already a competitive model from the same training pipeline.

\textbf{However, same-family distillation works well for function calling}: Qwen3-1.7B $\to$ Qwen3-4B improves full\_acc from 59.8\% to 73.7\% (+13.9pp), and Qwen3-4B $\to$ Qwen3-8B improves from 73.0\% to 76.8\% (+3.8pp).
This suggests that function calling capability has a different capacity profile from mathematical reasoning within the Qwen3 family.

\subsection{Cross-Family Function Calling (Full Table)}

\begin{table}[h]
\centering
\caption{Full function calling results across all model families (LoRA, best step).}
\small
\begin{tabular}{@{}lcccccc@{}}
\toprule
& \multicolumn{2}{c}{\textbf{Qwen}} & \multicolumn{2}{c}{\textbf{Gemma}} & \multicolumn{2}{c}{\textbf{Llama}} \\
\cmidrule(lr){2-3} \cmidrule(lr){4-5} \cmidrule(lr){6-7}
$N$ & Parse & Full & Parse & Full & Parse & Full \\
\midrule
50   & \textbf{98.3} & 57.2 & \textbf{60.0} & \textbf{30.9} & 99.3 & 44.0 \\
100  & 91.3 & \textbf{61.3} & 11.9 & 6.4  & 99.8 & 40.5 \\
150  & 92.5 & 61.5 & --- & ---  & 99.5 & \textbf{59.0} \\
200  & 90.2 & 54.5 & --- & ---  & 99.7 & 51.2 \\
Full & 86.7 & 58.2 & 11.7 & 3.9  & 99.5 & 32.0 \\
\bottomrule
\end{tabular}
\end{table}

Notable patterns: (i)~Gemma shows extreme sensitivity to position limit, with only $N=50$ producing usable results. (ii)~All Llama configurations achieve $>$99\% parse rate, indicating that JSON format is learned universally in the first 50 tokens. (iii)~Optimal position matches typical function call length for each tokenizer: Gemma $\sim$50, Qwen $\sim$100, Llama $\sim$150 tokens.

\subsection{Gemma Cross-Task Results}

\begin{table}[h]
\centering
\caption{Gemma-2-2B $\to$ Gemma-3-4B results across all tasks (LoRA pos-100, best step).}
\small
\begin{tabular}{@{}lccc@{}}
\toprule
\textbf{Task} & \textbf{Baseline} & \textbf{Best Distilled} & \textbf{$\Delta$} \\
\midrule
Math (avg@4) & 13.45\% & 27.20\% & +13.75pp \\
Funcall (full\_acc) & 73.50\% & 82.50\% & +9.00pp \\
HE+ (pass@1) & 20.1\% & 26.2\% & +6.1pp \\
MBPP+ (pass@1) & 34.7\% & 42.3\% & +7.6pp \\
\bottomrule
\end{tabular}
\end{table}

Cross-family positional distillation improves all four task metrics for the Gemma pair, despite requiring cross-tokenizer alignment (only 69\% vocabulary overlap by string matching).


\section{Position Is Not Reducible to High Entropy}
\label{app:position_vs_entropy}

A natural alternative to positional distillation is \emph{entropy-based} selection: pick the tokens where the student is most uncertain (highest surprise $-\log p_y$ on the sampled token), regardless of where they sit in the sequence. If head tokens were simply ``the high-entropy tokens that happen to come first,'' position would be a redundant proxy. We show empirically that this is not the case: position carries information beyond what entropy ranking provides.

We collected $74{,}549$ student-generated tokens (Qwen2.5-Math-1.5B on MATH-500) and computed per-token surprise $-\log p_y$. We then asked: among the first $N$ tokens, what fraction fall in each global surprise quartile? A baseline of $25\%$ in each quartile would mean the head is uncorrelated with surprise; full reduction to entropy would mean $100\%$ in Q1.

\begin{table}[h]
\centering
\caption{Surprise-quartile distribution among the first $N$ tokens. Baseline (uniform): $25\%$ per quartile. Q1 = top 25\% surprise globally, Q4 = bottom 25\%.}
\small
\begin{tabular}{@{}rrrrrrr@{}}
\toprule
\textbf{Head $N$} & \textbf{$n$} & \textbf{Q1 (top 25\%)} & \textbf{Q2} & \textbf{Q3} & \textbf{Q4 (bot 25\%)} & \textbf{Q1 enrichment} \\
\midrule
25  &  2{,}500 & \textbf{70.0\%} & 14.0\% &  8.6\% &  7.4\% & $2.8\times$ \\
50  &  5{,}000 & \textbf{61.0\%} & 14.7\% & 11.2\% & 13.1\% & $2.4\times$ \\
75  &  7{,}500 & 57.0\% & 15.7\% & 11.8\% & 15.6\% & $2.3\times$ \\
100 &  9{,}991 & 53.2\% & 17.1\% & 12.4\% & 17.3\% & $2.1\times$ \\
150 & 14{,}941 & 47.1\% & 18.9\% & 13.4\% & 20.5\% & $1.9\times$ \\
200 & 19{,}856 & 42.6\% & 20.0\% & 14.3\% & 23.1\% & $1.7\times$ \\
300 & 28{,}958 & 37.5\% & 21.1\% & 15.9\% & 25.5\% & $1.5\times$ \\
500 & 44{,}553 & 31.8\% & 22.3\% & 18.8\% & 27.1\% & $1.3\times$ \\
\bottomrule
\end{tabular}
\end{table}

\begin{table}[h]
\centering
\caption{Cumulative coverage: of all tokens above each global surprise threshold, what fraction lies in the first $N$ positions. \textbf{Reverse view}: of all $>p_{95}$ tokens, the median position is $172$, not $50$.}
\small
\begin{tabular}{@{}rrrrrr@{}}
\toprule
\textbf{Head $N$} & \textbf{\% of all tokens} & \textbf{cover $>p_{75}$} & \textbf{cover $>p_{90}$} & \textbf{cover $>p_{95}$} & \textbf{cover $\sum$ surprise} \\
\midrule
50   &  6.7\% & 16.4\% & 21.4\% & 25.1\% & 27.2\% \\
100  & 13.4\% & 28.5\% & 34.2\% & 38.4\% & 40.2\% \\
200  & 26.6\% & 45.4\% & 49.8\% & 53.4\% & 54.7\% \\
500  & 59.8\% & 76.1\% & 78.6\% & 79.8\% & 81.2\% \\
1000 & 98.5\% & 99.3\% & 99.5\% & 99.5\% & 99.5\% \\
\bottomrule
\end{tabular}
\end{table}

\textbf{Key findings:}
\begin{itemize}
\item Head tokens are \emph{enriched} in high surprise but not exclusively high surprise. The first 50 tokens are $2.4\times$ enriched in Q1 but still contain $13.1\%$ Q4 (low-surprise) tokens; the first 100 tokens contain $17.3\%$ Q4. Position and surprise are correlated but not collinear.
\item Head positions cover only a \emph{fraction} of high-surprise mass. Prefix-100 captures $38.4\%$ of all $>p_{95}$ surprise tokens; even prefix-500 only reaches $79.8\%$. To cover $95\%$ of $>p_{95}$ tokens one would need to keep the first $\sim$$846$ positions---essentially the entire response.
\item This rules out the hypothesis that prefix-100 wins by being a cheap entropy proxy. If high-entropy selection were the underlying mechanism, top-KL-100 (which by construction covers $>95\%$ of KL signal globally) should match or beat prefix-100. Instead, top-KL-100 underperforms prefix-100 by $7.3$pp avg@4 (\S\ref{sec:method}, Table~\ref{tab:n100}).
\end{itemize}

The most plausible interpretation is that \emph{position itself} is a load-bearing signal: an early-position low-surprise token can still be load-bearing for the trajectory through autoregressive cascade, while a late-position high-surprise token may already lie on an unrecoverable error path. Entropy ranks tokens by local uncertainty; position ranks them by causal influence on the rest of the generation. Empirically, the latter dominates for distillation utility.


\end{document}
